\chapter{Discussion}
\label{chap:discussion}

\section{Interpretation of results}
\label{sec:disc_interpretation}

The results of Chapter~\ref{chap:results} do not identify a single best inverse-kinematics method. Each method tested dominated some evaluation regime on some dataset, and each was beaten elsewhere. This section reads the results as a coherent story rather than as isolated experiments, argues that the regimes depend on a small number of physical properties of the system, and proposes explanations for the least obvious findings.

\paragraph{Every classical baseline failed, but for different reasons.}
The single-segment PCC of Section~\ref{sec:res_pcc_single} reached a $22$\,mm mean forward residual on the reachable subset of Dataset 2, well above the $\sim$$2$--$10$\,mm scale that would make it a competitive baseline on a workspace roughly $100$\,mm across. The diagnostic analysis attributed this to a dominant elongation signal ($r = +0.737$ between total pressurisation and tip norm) that a constant-arc-length PCC model has no degree of freedom to represent. The two-segment PCC of Section~\ref{sec:res_pcc_multi} reached geometric inverse residuals of $3 \times 10^{-3}$\,mm---essentially zero---yet its closed-loop error on Dataset 1 was $51$\,mm, roughly four times larger than the worst data-driven baseline. Here the failure is different: the geometry is fine, but the multi-valued two-segment inverse hands the linear pressure-calibration step a set of training frames whose chamber lengths collide on very different pressure commands, and ordinary least squares recovers the per-chamber mean pressure. These two failure modes are instructive because they show that a correct geometric model does not guarantee a competitive inverse: the inverse has to be well-posed both in configuration space and in actuation space, and neither dataset provides that joint well-posedness for a pure PCC pipeline.

\paragraph{Data-driven regression lets the inverse be as ill-posed as the dataset allows.}
Direct-regression networks do not need the inverse to be well-posed: when the inverse maps one tip to many pressures, the least-squares loss gently pulls the network toward the training-distribution average pressure for that tip. This is visible on Dataset 1, where the tip-only feed-forward network reached $R^{2} = 0.32$---a value that reflects a genuinely under-determined inverse, not a training failure, as confirmed by the early validation plateau at epoch $4$ (Section~\ref{sec:res_mlp_tip}). On Dataset 2, where three valves map into three tip coordinates, the inverse is essentially single-valued on the observed workspace, and the same architecture reached $R^{2} = 0.93$. This contrast is one direct argument for diagnosing the task by actuation redundancy rather than by architecture choice: the direct-regression target is the per-tip posterior mean over valid pressures, and its predictive ceiling is set by the concentration of that posterior, not by the expressive power of the network.

\paragraph{Mid-marker input lifts the posterior concentration; temporal context amplifies it.}
Adding the mid-marker coordinates to the Dataset 1 input was the largest single-factor improvement reported in this thesis ($R^{2}$ from $0.32$ to $0.44$ on the feed-forward network, from $0.28$ to $0.51$ on the recurrent network; Sections~\ref{sec:res_mlp_tipmid}, \ref{sec:res_lstm}). The information-theoretic reading is that the tip under-identifies the internal configuration of a two-segment arm, and that the mid marker supplies a partial observation of that configuration, narrowing the set of pressures consistent with the observation. Temporal context alone does \emph{not} produce this effect---the tip-only recurrent network is marginally worse than the tip-only feed-forward network on Dataset 1---yet the combination of mid-marker input and temporal context is better than either alone. The likely reason is that the recurrent network extracts the trajectory of the internal configuration rather than the trajectory of the tip, and that trajectory is a meaningful feature only once the internal configuration is itself observable. This is consistent with the hysteresis hypothesis that motivated Experiment~3 but narrows it: history helps when the features it has access to are informative enough for the history to carry physical meaning, and does not help when those features under-identify the system.

\paragraph{The Fang 2022 trade-off is a property of multi-valued inverses, not of neural architectures.}
The Jacobian-learned iteration of Section~\ref{sec:res_jacobian} was predicted to sacrifice pressure-space accuracy for position-space accuracy on a multi-valued inverse, on the grounds that direct regression averages valid solutions while a Jacobian iteration commits to one \cite{fang2022jacobian}. The measured trade-off on Dataset 2 was a $33\%$ increase in pressure MAE against a $71\%$ reduction in closed-loop position error (Table~\ref{tab:res_jacobian_pos}), and on Dataset 1 the trade-off is similar in sign but smaller in magnitude. The same reversal would be expected for any pair of methods that chooses posterior-mean versus posterior-mode behaviour on a multi-valued inverse, regardless of their implementation details. The practical consequence is that Jacobian iteration and direct regression are not ranked on a single ``quality'' scale but instead measure different objectives: a controller that places the tip in a desired position should adopt a Jacobian iteration; a system that reports pressure MAPE to a downstream consumer should adopt direct regression.

\paragraph{Sim-to-real transfer is dominated by workspace geometry, not by sample count.}
The negative sim-to-real result on Dataset 2 (Section~\ref{sec:res_sim_transfer}) is worth reading with care. Strategy~B---sim pre-training followed by real fine-tuning---matched Strategy~A at every real-data budget tested, including the $10\%$ budget at which one would most expect a prior to help. Strategy~C, which used only synthetic data, failed with $R^{2} = -1.71$ on Dataset 2. The synthetic workspace produced by the calibrated rod of Section~\ref{sec:res_sim_cal} was five times the linear extent of Dataset 2 on every axis and centred in a different region of tip space, so the pre-training distribution had almost no overlap with the test distribution. Under these conditions the pre-training is not a prior toward the right answer but a prior toward an irrelevant answer, and fine-tuning on a small fraction of real data discards most of the pre-trained weights. The more general reading is that a Manti-inspired rod is not a drop-in prior for an undocumented robot: the value of any such prior is bounded above by the overlap between its synthetic workspace and the real operating region, and the $720$-sample count that limited the present simulation is not the reason this ceiling was hit.

\paragraph{A compact decision criterion.}
Across the six methods and two datasets tested here, the right inverse-kinematics method for a new pneumatic soft-robot arm is determined less by the method itself than by three properties of the arm and its intended use. First, actuation redundancy: if there are more actuation degrees of freedom than tip degrees of freedom, a Jacobian-based iteration can be expected to drive closed-loop error to the tolerance floor, as observed on the six-chamber Dataset 1 arm. Second, sensing coverage: if the tip under-identifies the internal configuration, additional observations---mid markers, full-shape sensing, or any other partial observation of the internal state---lift the predictive ceiling more than any network-architecture change, as observed on Dataset 1 with the mid-marker ablation. Third, target-domain overlap: a simulator prior is useful only when its synthetic workspace covers the intended operating region of the real arm; otherwise the most data-efficient route is still to train on real data alone, as observed on Dataset 2. These three properties together recover most of the regime-specific findings of Chapter~\ref{chap:results} and give a compact decision criterion that does not privilege any single method.

\section{Comparison with Published Baselines}
% TODO: explicit comparison with Ma et al. (2022) (MAPE~<5\%) and other
%       benchmarks from docs/references/giorelli_nn_vs_jacobian.md.

\section{The Role of Mid-Segment Information (Dataset 1)}
% TODO: discuss why tip-only IK is underdetermined for two-segment arms
%       and what this implies for sensing requirements.

\section{Temporal Modelling and Hysteresis}
% TODO: when does LSTM help, when not; limits of the tested sequence
%       lengths.

\section{Jacobian vs. Direct Regression}
% TODO: generalisation, sample efficiency, robustness to out-of-distribution
%       targets.

\section{Sim-to-Real Observations}
% TODO: what transfers, what does not; the reality gap for pneumatic
%       actuation dynamics.

\section{Limitations}

\subsection{No Physical Robot}
% TODO: all evaluation is on pre-recorded datasets; no closed-loop testing.

\subsection{Dataset Coverage}
% TODO: two datasets, two geometries; broader generalisation unverified.

\subsection{Simulator Fidelity}
% TODO: PyElastica limitations (no air dynamics, simplified hysteresis).

\subsection{Metric Choice}
% TODO: tip-only error ignores shape; implications for downstream use.

\section{Implications}
% TODO: practical take-aways for practitioners choosing an IK strategy
%       for a new pneumatic soft arm.
